\providecommand{\devnum}[1]{#1}

\section{Limitations and Threats to Validity}
\label{sec:limits}

\paragraph{The shared pool is our design.}
CodeBench gives each logged-in student a private container, so on the real
platform one job never waits behind another and there is no shared
judging queue \cite{codebench_dataset}. We use its logs as a demand trace: the
arrival instants and the per-job costs are recorded; the server pool, the
queue and the scheduler are ours. Section~\ref{sec:data} sets out what supports
that design, that a fixed pool of judge hosts on one queue is a standard
deployment \cite{domjudge_manual,peveler2019sandbox} and that pooling replaces
\devnum{1{,}636} concurrent containers with \devnum{14} servers on the same
demand \cite{smith1981pooling}. What it cannot support is any statement about
waiting on the real platform, the pooled side being simulated throughout. The
container counts belong to the 44-fold overlay and to no deployment, and
a container also serves the student's editing session, so the ratio is not a
cost saving.

\paragraph{Open-loop replay.}
Arrivals come from the log at their recorded instants and do not react to our
scheduler. Feitelson states the objection at its strongest, that
``logged workloads actually contain a `signature' of the logged system''
\cite{feitelson2021resampling}; a trace-driven simulation issues requests
irrespective of system state, so a scheduler is neither rewarded with the work
it would attract nor punished by the users it would drive away
\cite{shmueli2009simulation}. We cannot restore that feedback, our logs
carrying no session boundaries, no think times and no resubmission semantics,
but its direction is known: a closed loop returns work to a service that
answers quickly, so the omission suppresses the benefit of the policies that
shorten waits, and our reported gains are a lower bound. Replaying into a
differently sized pool is established practice
\cite{harcholbalter2003sizebased}, and the width coupling that distorts
parallel-job replays does not arise here, a job occupying one server whatever
$k$ is. One choice is deliberate. We raise load by superposing whole
class-semesters aligned on weekday and clock time; compressing inter-arrival
times, Feitelson notes, shrinks the daily load cycle along
with the gaps \cite{feitelson2021resampling}.

\paragraph{The guarantee is denominated in $L$.}
Every constant in Section~\ref{sec:theory} is a multiple of the per-job time
limit. The additive constant $(3-2/k)L$ is an exact supremum and no
artefact of the proof, it grows with $k$ from $L$ at a single server towards
$3L$, and a wrapper of this family cannot promise less than $L$, whatever base policy
it is given: an excess of at least $L$ is forced on some job whenever the base policy
reads a service time only at completion, so no mechanism of this family can promise
less than $\Omega(L)$.
The obvious escape from the $\Omega(L)$ floor is kill-and-restart tiering, and it
fails. Kill-and-restart tiering runs every
job for a short first quantum $\tau$ and restarts the survivors, discarding the
killed attempt; on $n$ jobs of length $L$ arriving together it leaves some job
with an excess of $n\tau/k - L$ over first-come first-served, and that excess
is wasted work rather than overtaking, so no budget removes it. The light-job
constant does not improve either, a second-tier piece being still a
non-preemptive block of length $L$: holding $\tau$ fixed and varying $L$, the
worst light-job additive constant we measure grows as about \devnum{1.1}$L$ and
is flat in $\tau$. The mechanism is TAGS \cite{harcholbalter2002tags} with the
dedicated hosts removed, and removing them is what breaks it: in TAGS the kill is a
routing device and the work is redone on the next host, whereas in a judging queue the
kill is terminal.

\paragraph{Two applicability conditions, each with traces that fail it.}
The first is that $L$ be comparable to the waits that matter, and four of the six
workloads fail it: the first-come first-served p99 wait is \devnum{0.021},
\devnum{0.136}, \devnum{0.373} and \devnum{1.78} times $L$ on ACcoding, the compute
farm and the two continuous-integration pools, against the floor of about
\devnum{3}$L$ below which no promise can be stated. On the Intel
Netbatch compute farm \cite{netbatch2012trace,shai2014netbatch} the service cap
is \devnum{5.7} days while the first-come first-served p99 wait in the worst
window we build is \devnum{18.7} hours, so the smallest expressible promise is
about \devnum{17} days and the guard never fires at any admissible budget;
predicted-cost ordering also makes the p99 wait worse than first-come
first-served there under sustained overload. The guard results we report on the
four say what the ordering mechanism does there; none of them reports a promise.
The second condition is that the cost
distribution be heavy-tailed. On the two Firefox CI hardware pools
\cite{taskcluster_docs,treeherder_docs} the longest \devnum{1\%} of jobs carry
\devnum{3.4\%} and \devnum{5.2\%} of the work, at coefficients of variation
\devnum{0.61} and \devnum{0.87}, and there even true-size ordering makes the
p99 wait worse than first-come first-served on three of the four windows we
ran. Section~\ref{sec:exp_cross} gives both conditions in numbers.

\paragraph{The guarantee is per job, not per quantile.}
The bound holds for each job separately and says nothing about a percentile. A
guarded run can have a worse deadline-window p99 than the baseline its guarantee
is stated against, by as much as the promise itself, and Table~\ref{tab:adv}
shows it doing so under two of the three corrupted scores at the busiest load
while the worst single job stays inside the promise.

\paragraph{Development data, and what is still sealed.}
Every CodeBench number in this paper is development data. The guard family and
an early parameter grid were explored on the same overlay the tables report,
before the three pre-stated grids were searched on validation overlays that
exclude the evaluation semester, so the selection is validation-only inside a
family that is not, and no number here is out of sample. The rule is a
worst-cell rule, so it is sensitive to which validation cells it sees: at the
middle promise, dropping the one cell that decides the choice moves it to a
different point of the same grid, and five overlays are week-shifts of the
same class-terms, so they are not five independent tests. We do not describe any of
these parameters as validated. The guarantee does not depend
on the parameters being well chosen; only the benefit does. The sealed semesters
are opened once, after the method, the parameters and the primary metric have
been frozen, and their numbers are reported separately.

\paragraph{The budget shape is selected on one set of cells, and does not always
transfer.} The promise holds for any budget under the cap, so nothing about the
guarantee depends on the shape. What the shape decides is how much of the base
policy's benefit survives, and that we select on \devnum{15} multi-server
validation cells; those cells are the whole evidence we have for it. The
single-server configuration shows the limit of that evidence. At $G = \devnum{1200}$ s the
queue-length shape wins the joint comparison on those cells and closes
\devnum{0.701} of the gap at $k = 1$, where the age-relative shape closes
\devnum{0.895}; at the two tighter promises the same choice does transfer.
We report the number the pre-stated rule produces and do not re-tune on the
configuration that exposes it, and an operator whose pool is much smaller than
the validation pool should re-run the selection rather than take our shape. The
three shapes also differ in where the harm falls, not only in how much
gap they keep, so which of them is preferable depends on a judgement about jobs
first-come first-served would have started at once, and that judgement is the
operator's.

\paragraph{What the recorded costs mean.}
The cost we replay is the platform's recorded execution time, a wall-clock
duration for one evaluation that includes interpreter start-up. Runs of the
same student that overlap stretch one another, and between \devnum{38\%} and
\devnum{79\%} of jobs longer than a second end in the same second as
another block of that student, so part of the recorded cost is contention we
then re-impose inside the pool. How the platform enforces its
\devnum{60}-second limit is documented nowhere we could find. We wrote to the
maintainers about both points and had no answer.
ACcoding \cite{accoding2024} records no submission timestamps, so its arrival
process is synthetic. The Firefox CI check, the one trace whose queue waits are
themselves recorded, fits an effective capacity and a setup term, both refitted
on one chronological part of each window and tested on the other; out of sample
the mean wait, the per-hour profile and the per-job ordering survive and the
90th percentile does not, landing up to \devnum{21\%} from the recorded one and
systematically low on the larger pool, so the simulated distribution is too steep
in its upper middle. The split is chronological inside one \devnum{21}-day and
one \devnum{7}-day window of two pools, not a different cluster, quarter or
workload, and capacity is held constant over the test part although the real
machine count varies.

\paragraph{Not studied.}
Preemption is outside the model, so nothing here applies to a system that can
checkpoint or suspend. The wrapper applies unchanged to a learned or
reinforcement-learned base policy, its bound assuming nothing about the base,
but we did not test one. Capacity is fixed throughout, so this is not a
capacity-planning study. The reported runs compare policies at a common set of
enqueue instants and do not charge the cost of computing features and running the
predictor. Section~\ref{sec:prediction} measures that cost and charges it in a
sensitivity, where it moves the closed gap by less than \devnum{0.001} at every
load; what is missing is not the latency but the rest of an end-to-end system.
Nothing here was deployed. A deployed serverless scheduler of the kind this
wrapper is meant for already exists: ALPS learns a per-function rank and a
time-slice cap from an offline shortest-remaining-processing-time replay and
leans on the Linux completely fair scheduler to keep any one invocation from
starving \cite{fu2024alps}, which is a mechanism without a proved per-job bound.

% The Conclusions section that follows was the closing \paragraph{Conclusion.} of this
% section until the sixth revision pass; MDPI's structure expects a numbered Conclusions
% section and the referee asked for one (RR-18).  It is kept in this file rather than
% given its own sections/ file so that main.tex's list of \papersection calls is untouched.
\section{Conclusions}
\label{sec:conclusions}

We asked whether a non-preemptive shared execution service can be reordered by
predicted cost without exposing any single job to an unbounded wait. What the
evidence supports is that work is the currency the excess is written in. A
job's delay relative to first-come first-served equals its net
overtaken work divided by the number of servers, up to a constant multiple of
the time limit, and a budget on that quantity is sufficient and necessary for a
per-job guarantee. Counting overtakes is not necessary for such a guarantee, and
buys one only at the worst-case charge of the time limit per pass, which on our
traces overstates the work that actually passes by one to two orders of
magnitude. The wrapper built on it holds every job to a promise the
operator states in seconds, whatever the predictor does, at a cost of about an
eighth of the gap predicted-cost ordering closes at the busiest load we simulate
and at a promise of ten time limits. The method needs a time limit comparable to
the waits at issue and a heavy-tailed cost distribution, and four of our six
workloads fail the first of those.

Two questions are left open. The identity is stated in the dispatch-sequence
currency, where the constant is exactly determined in both directions; under the
executed-work convention the upward constant is unchanged and the downward one is
known only to lie in $((k-1)L, 2(k-1)L]$, and closing it would make the two
conventions interchangeable. And Algorithm~\ref{alg:guard} is not claimed to be the
least restrictive rule carrying its guarantee: whether some wrapper overrides the
base policy strictly less often and still promises $G$ is a question we have not
settled in either direction.
